% Original friend-tree dynamics, not a scalar independent-family replacement.
% Dependencies: the aggregate disagreement construction and m(s)<=19/20
% for s>=11/20, the exterior killed-reservoir survival lemma, and the
% established inert pressure r_infinity(mu)=1.

\section{Critical boundary values and core-comparable eigenfunctions}
\label{cr:section}

The inputs from the original manuscript\cite{Long} are its aggregate
disagreement coupling, the certified infection mass at most $19/20$ for
discounts at least $11/20$, the exterior reservoir survival probability
at least $1/6$, the uniform supercritical resolvent bound, and the inert
pressure equation $r_\infty(\mu)=1$. The proof below strengthens those
inputs to fixed-seed critical statements.

Put $s_0=11/20$ and $q=19/20$.  Let $\mathcal X^\partial$ be the
class of finite newborn environments containing at least one inert
stub.  The whole environment, including every stub, is retained when
retagging, so this class is invariant under the genealogical kernel.
For $x\in\mathcal X^\partial$, set
\[
 V_{s,T}^{\rm tag}(x)=\int_0^T e^{-st}F_x(t)\,dt,
 \qquad W_T^\partial(s)=\sup_{x\in\mathcal X^\partial}
 V_{s,T}^{\rm tag}(x).
\]
Use $V_s^{\rm tag},W^\partial(s)$ for infinite horizons, and abbreviate
$V_s^\infty=V_s^{\rm tag}(x_\infty)$.  The already proved uniform
resolvent theorem gives $W^\partial(s)<\infty$ for $s>\mu$.
For finite $T$, Yule domination gives finiteness for every $s>0$.

\subsection{A directional version of the disagreement gradient}

Consider a finite retained core tree with any finite number of stubs.
Mark an old core $u$ and count only its retained descendant component,
the side of its parent edge containing $u$, or the whole component if
that parent is an inert stub.  Its present cores and all their future
descendants are included in the reward.  Let $v$ be an old core
outside this counted component.  Write its old path to $u$ as
\[
 v=v_r,v_{r-1},\ldots,v_1,v_0=u,
\]
and let $d_i\ge2$ be the full degree of $v_i$ in the state before an
additional leaf is attached at $v$.  The path contains only real
cores and persists forever.

\begin{lemma}[Directional root-leaf gradient]
\label{cr:directional}
If the retained environment has at least one stub, then for $s\ge s_0$
and finite $T$ the difference in the two tagged discounted rewards is
at most
\begin{equation}
 \frac{2W_T^\partial(s)}{s(1-q)}
 q^r\prod_{i=1}^r\frac1{d_i-1}.
 \label{cr:directional-bound}
\end{equation}
For an additional leaf at $u$ itself, use $r=0$ and the empty product
one.  The estimate is uniform in the old state and its finite stubs.
\end{lemma}
\begin{proof}
Use the aggregate leaf-gain and promotion coupling from the existing
disagreement construction.  Start the infection at $v$.  At the first
infection of any other eligible core, its two degrees differ by one;
freeze its then-present common neighbors, and charge every new core
family born from it subsequently.  The infected old eligible set
remains connected.  A charged branch can contact the old tree only
at the core from which it was born.  It therefore cannot bypass an
edge of the displayed old path.  Its repeated feedback is retained
in the coupling, rather than deleted.

At a noninitial infected core of smaller full degree $d$, use the
existing reservoir $Q=Z+R+B$: degree difference $Z$, union promotions
$R$, and infections of frozen outward neighbors $B$.  The existing
argument includes all repeated feedback through the predecessor,
infected outward neighbors and charged branches, and gives upward
rate at most $Q+5/2$.  At the initial infection there is no infected
predecessor and there may be $d$, rather than $d-1$, frozen outward
neighbors.  In this case
\[
 |a^+-a^-|\le Z+R+B/2,
 \qquad \text{rate}(B\text{ increases})
 \le \frac{dZ}{d(d+Z)}\le1.
\]
Thus the reservoir's upward rate is at most $Q+2$, which is also
bounded by $Q+5/2$.  In both cases it is dominated, conditionally on
the entire infection history, by the same pure-birth chain $Y$ with
$Y(0)=1$ and rate $n+5/2$ at $n$.

For a specified still-uninfected frozen neighbor, its infection
intensity is at most $Z/[d(d+Z)]$.  Its expected discounted first
infection indicator is consequently at most
\[
 e_d(s):=\mathbb E\int_0^\infty
 e^{-st}\frac{Y(t)}{d(d+Y(t))}\,dt
 =\frac{m_d(s)}{d-1}\le\frac q{d-1}.
\]
This is a bound for each specified edge, not merely for the sum of
the offspring.  It remains valid for the initial infection with $d$
candidates, because the reservoir domination just verified is the
same.  Full degrees never decrease, so the smaller degree upon
infection at $v_i$ is at least its initial degree $d_i$.

The first infection of $u$ must traverse each old path edge in order.
Conditional expectation at the successive infections therefore gives
\[
 \mathbb E[e^{-s\tau_u};\tau_u<\infty]
 \le q^r\prod_{i=1}^r(d_i-1)^{-1}.
\]
No independence between path vertices or charged families is used.
Before $\tau_u$, both counted descendant histories agree.  After it,
run the existing disagreement exploration restricted to the counted
component.  Its expected discounted number of eligible infections,
relative to $\tau_u$, is at most $1/(1-q)$.  At each such infected
core the union-promotion rate is at most one; each promotion charges
at most two newborn families, each with remaining discounted reward
at most $W_T^\partial(s)$.  The whole retained environment of every
family still contains an inert stub.  Its expected discounted charge
is at most $2W_T^\partial(s)/s$.  This proves
\eqref{cr:directional-bound} exactly as in the original gradient proof.
Localization at population and infection generation, followed by
Tonelli, justifies the estimates at finite horizon.
\end{proof}

\subsection{Uniform relative approximation by finitely many profiles}

For a newborn $x$, use the existing cut map $\pi_{R,M}$: traverse the
initial core from the tag, cut a branch at its first core of initial
full degree at least $M$, or at depth $R$, and replace each cut edge
by an inert stub.  Initial leaf counts and full core-neighbor counts
of retained cores are preserved.  All retained full degrees are less
than $M$, and retained depths are less than $R$, so the range is
finite.  An original stub either remains retained or lies beyond a
cut edge, which becomes a stub.  Consequently
\[
 x\in\mathcal X^\partial\quad\Longrightarrow\quad
 \pi_{R,M}x\in\mathcal X^\partial.
\]

\begin{lemma}[Relative finite-profile estimate]
\label{cr:profiles}
For $s\ge s_0$, $R\ge1$, $M\ge3$, and finite $T$,
\begin{equation}
 \left|V_{s,T}^{\rm tag}(x)
       -V_{s,T}^{\rm tag}(\pi_{R,M}x)\right|
 \le\kappa_{R,M}(s)W_T^\partial(s),
 \label{cr:relative-profile}
\end{equation}
uniformly over $x\in\mathcal X^\partial$, where
\[
 \kappa_{R,M}(s)=\frac{2}{s^2(1-q)}
 \left(\frac1M+\frac12q^{R-1}\right).
\]
\end{lemma}
\begin{proof}
Keep the exterior evolution, but delete the incoming arrow channel
across every chosen cut edge $w\to v$, with $v$ retained and $w$
removed.  The retained projection of this modified process follows
exactly the inert profile law.  Outgoing choices $v\to w$ alter only
the exterior, and are null choices in that projection.  No later
branch can create another connection across this initial tree cut.

The actual retained process differs from the profile generator only
by predictable leaf inputs $q_{wv}(t)\Delta_v$, where
\[
 0\le q_{wv}(t)=1/D_w(t)\le1/D_w(0).
\]
Apply the finite-horizon discounted backward Dynkin identity to the
tagged profile value function, evaluated along the actual retained
state.  Each perturbation is bounded by Lemma~\ref{cr:directional}.
The counted component may already contain many old cores at the
time of perturbation; this is permitted in that lemma.  The old path
from $v$ to the tag persists, and every current full degree on it is
at least its initial core-neighbor count $c_i$, including stubs.
For an internal vertex on a cut path, $c_i\ge2$.

The newborn tag has initial core-neighbor count one.  Consider a
nonbacktracking walk on the initial core-and-stub tree: the first
step is forced, and at each later core $i$ it chooses uniformly
among the $c_i-1$ forward core neighbors or stubs.  For the path
ending with the cut edge $w\to v$, with $v$ at depth $r$, its cylinder
probability is
\[
 p(w)=\prod_{i=1}^{r}(c_i-1)^{-1}.
\]
For a cut adjacent to the tag this is the empty product one.  First-cut
paths are prefix-free, so $\sum_w p(w)\le1$.  The directional gradient
contributes at most $q^rp(w)$, multiplied by its common prefactor.

For cuts caused by high degree, $1/D_w(0)\le1/M$ and $q^r\le1$.
For all remaining depth cuts, $r=R-1$ and
$1/D_w(0)\le1/2$.  Summing these two disjoint collections gives
\[
 \sum_{w\to v}\frac{q^{\operatorname{depth}(v)}p(w)}{D_w(0)}
 \le\frac1M+\frac12q^{R-1}.
\]
Finally $W_{T-t}^\partial(s)\le W_T^\partial(s)$ and
$\int_0^T e^{-st}dt\le1/s$.  These estimates yield
\eqref{cr:relative-profile}.  The cut set is finite for each initial
environment.  Its size never enters the bound, because the path
probabilities have already been summed.  Population localization and
finite-time moments justify Dynkin's identity before taking the limit.
\end{proof}

For example, take $R=161$ and $M=1024$.  The rational inequality
$2\cdot19^{16}<20^{16}$ implies $q^{160}<1/1024$, and therefore
\[
 \kappa_{161,1024}(s)
 \le\frac{40}{(11/20)^2}\frac3{2048}
 =\frac{375}{1936}<\frac14
 \qquad(s\ge s_0).
\]
No numerical approximation is needed for this choice.

\subsection{Each stub profile is dominated by a fixed inert seed}

\begin{lemma}[Permanent guardian embedding]
\label{cr:guardian}
For each fixed finite newborn profile $p\in\mathcal X^\partial$,
there are $h_p<\infty$ and $a_p>0$ such that
\begin{equation}
 F_\infty(h_p+t)\ge a_pF_p(t)\qquad(t\ge0).
 \label{cr:guardian-time}
\end{equation}
Consequently, on every bounded interval $0<s\le S$,
\[
 V_s^{\rm tag}(p)\le a_p^{-1}e^{Sh_p}V_s^\infty.
\]
\end{lemma}
\begin{proof}
Let $b\ge1$ be the number of profile stubs.  Replace every stub by
an edge to a separate new real guardian with a private real leaf.
Choose one guardian as the original inert root and give that guardian
the one original inert stub.  The augmented real graph is a finite
tree with exactly one inert stub, and the guardian root has its
private leaf.  It can be obtained from $x_\infty$ by finitely many
specified leaf attachments, since every real target has a real
neighbor with positive rate of choosing it.  The profile tag is a
leaf of its initial core skeleton with its own sole leaf; build its
promotion last.  Requiring the finitely many prescribed clocks and
choices in bounded time, with other clocks silent during preparation,
gives an event of probability $p_p>0$ on which the augmented tree is
exact at a deterministic time $h_p$.  All its cores descend from the
original inert root.

Restart independent drivers at $h_p$.  For each guardian, use exactly
the exterior killed-reservoir construction from the uniform boundary
removal lemma.  The reservoir $X$ starts at one, births at rate $X$
using that guardian's independent leaf-birth Poisson measure, and
loses one at every guardian source-clock ring while positive.  At
such a ring kill the reservoir if its independent uniform mark
$U\le1/X$.  The existing explicit bounded subharmonic function proves
that its permanent survival probability is at least $1/6$.

Couple $X$ below the actual guardian leaf count.  Actual leaf births
dominate the reservoir births, a source ring removes at most one
actual leaf, and all immigration from the profile or other exterior
branches only adds leaves.  Thus $D_g\ge L_g\ge X$ before killing.
Choose the guardian's unique interior neighbor exactly when
$U\le1/D_g$.  This is a correct uniform-neighbor construction.  On
reservoir survival, $U>1/X\ge1/D_g$ at every guardian ring, so the
guardian never sends an arrival to the profile.

These survival events use only the independent guardian driver arrays,
not the interior driver array.  They allow all guardian promotions
and all exterior side-branch evolution; they do not require a guardian
to remain a star.  Their joint probability is at least $6^{-b}$.
On their intersection the interior has exactly the prescribed profile
law, using unchanged independent interior drivers.  Outward interior
choices only add leaves at guardians, which cannot spoil the lower
reservoir coupling.  Conditional on the whole preparation history,
the expected tagged profile reward times the intersection indicator
is therefore at least $6^{-b}F_p(t)$.

That tagged reward is part of the original inert root's total core
reward.  Taking $a_p=p_p6^{-b}$ proves
\eqref{cr:guardian-time}.  Integrating from $h_p$ to infinity proves
the resolvent assertion.  This exact construction uses the assumption
that the profile has a stub; no such embedding is asserted for a
stub-free profile.
\end{proof}

\subsection{Critical divergence and a bounded positive eigenfunction}

\begin{theorem}[Uniform comparison to the inert resolvent]
\label{cr:uniform-inert}
There is a finite constant $K$ such that
\begin{equation}
 W^\partial(s)\le K V_s^\infty
 \qquad(\mu<s\le1).
 \label{cr:uniform-inert-bound}
\end{equation}
Moreover
\begin{equation}
 \boxed{V_\mu^\infty=\infty,
 \qquad\sum_{m\ge0}A_m^\infty(\mu)=\infty.}
 \label{cr:critical-divergence}
\end{equation}
\end{theorem}
\begin{proof}
Fix $R=161,M=1024$ and let $\mathcal P$ be the finite range of the
profile map on $\mathcal X^\partial$.  Lemma~\ref{cr:profiles} gives
at each finite horizon
\[
 W_T^\partial(s)\le
 \max_{p\in\mathcal P}V_{s,T}^{\rm tag}(p)
 +\frac{375}{1936}W_T^\partial(s).
\]
By Lemma~\ref{cr:guardian}, the maximum of the infinite-horizon values
is at most $K_0V_s^\infty$, where
$K_0=\max_{p\in\mathcal P}a_p^{-1}e^{h_p}<\infty$ for $s\le1$.
Absorption and then $T\to\infty$ prove
\eqref{cr:uniform-inert-bound} with
$K=K_0/(1-375/1936)$.

Suppose $V_\mu^\infty$ were finite.  The same finite-horizon argument
at $s=\mu>s_0$ would give $W^\partial(\mu)<\infty$.  Tonelli and the
first genealogical decomposition then give the bounded function
\[
 f(x)=\mu V_\mu^{\rm tag}(x)
      =1+\sum_{m\ge1}A_m(\mu;x),
 \quad 1\le f\le\mu W^\partial(\mu),
 \quad\mathcal K_\mu f=f-1.
\]
Let $B=\mu W^\partial(\mu)$.  Then $B>1$ and
$\mathcal K_\mu f\le(1-1/B)f$.  Positivity gives
\[
 A_m^\infty(\mu)
 \le (\mathcal K_\mu^m f)(x_\infty)
 \le (1-1/B)^m f(x_\infty).
\]
This contradicts the already established pressure equation
$\lim_m A_m^\infty(\mu)^{1/m}=1$.  Hence
$V_\mu^\infty=\infty$, and the genealogical resolvent identity proves
the second assertion.
\end{proof}

\begin{theorem}[A bounded positive critical genealogical eigenfunction]
\label{cr:eigenfunction}
On $\mathcal X^\partial$ there is a bounded function $h$ such that
\begin{equation}
 h(x_\infty)=1,\qquad
 \gamma_1\le h(x)\le K,\qquad
 \mathcal K_\mu h=h,
 \quad\gamma_1=\frac1{48}.
 \label{cr:eigenfunction-bounds}
\end{equation}
In particular, for every $m\ge0$,
\begin{equation}
 \boxed{\frac1K\le A_m^\infty(\mu)\le\frac1{\gamma_1}=48.}
 \label{cr:critical-lower-coefficients}
\end{equation}
The previous sharper upper bound $A_m^\infty(\mu)<43$ remains valid.
\end{theorem}
\begin{proof}
For $\mu<s\le1$, normalize
$h_s(x)=V_s^{\rm tag}(x)/V_s^\infty$.  The uniform comparison gives
$h_s\le K$.  The original minorization and the genealogical resolvent
identity give
\[
 V_s^{\rm tag}(x)
 =\frac1s\left(1+\sum_{m\ge1}A_m(s;x)\right)
 \ge\gamma_s V_s^\infty\ge\gamma_1 V_s^\infty.
\]
The first genealogical decomposition gives the exact equation
\[
 h_s=\frac1{sV_s^\infty}+\mathcal K_s h_s.
\]
Take a sequence $s\downarrow\mu$ and a diagonal subsequence on the
countable space $\mathcal X^\partial$.  Its pointwise limit $h$ has
the displayed bounds and normalization.  Monotone convergence gives
$V_s^\infty\to V_\mu^\infty=\infty$.  In the kernel expression,
$|h_s|\le K$ and
\[
 \mathbb E_x\sum_{v\text{ child}}e^{-\mu\sigma_v}\le1/\mu.
\]
Dominated convergence therefore gives
$\mathcal K_s h_s\to\mathcal K_\mu h$ at each state.  Thus
$\mathcal K_\mu h=h$.  Iteration and positivity yield
\[
 \gamma_1 A_m^\infty(\mu)
 \le(\mathcal K_\mu^m h)(x_\infty)=1
 \le K A_m^\infty(\mu),
\]
which proves \eqref{cr:critical-lower-coefficients}.
\end{proof}

The uniform profile comparison does not make the whole-environment
operator compact on $\ell^\infty$.  It is a relative approximation
for this particular family of resolvent rewards.  Accordingly, the
critical eigenfunction is compatible with the previously proved
critical essential spectral radius one.

\begin{corollary}[A genuine critical renewal boundary value]
\label{cr:renewal-boundary}
Let $F_s(z)=\sum_{m\ge0}A_m^\infty(s)z^m$ and define the signed
renewal germ $B_s(z)=1-1/F_s(z)$ near zero. Its Taylor coefficients are
the finite signed expressions
\[
 b_m(s)=
 \sum_{j=1}^{m}(-1)^{j-1}
 \sum_{\substack{n_1+\cdots+n_j=m\\n_i\ge1}}
 \prod_{i=1}^{j} A_{n_i}^\infty(s).
\]
Along the real interval
$0\le z<1$, $F_\mu(z)>0$ and this formula continues the germ, and
\[
 \frac{1-z}{48}\le1-B_\mu(z)\le K(1-z),
 \qquad \lim_{z\uparrow1}B_\mu(z)=1.
\]
\end{corollary}
\begin{proof}
Expand the reciprocal formally at zero to obtain the coefficient formula;
the kernel norm bound makes it an analytic germ, for example on
$|z|<s/2$. The bounded positive critical coefficients give
$1/[K(1-z)]\le F_\mu(z)\le48/(1-z)$ for real $z<1$.
Taking reciprocals proves both assertions.
\end{proof}

This boundary statement is a continuation along the positive real
axis.  It does not assert that the signed Taylor series for $B_\mu$
has radius at least one, nor that its coefficients are nonnegative.
Complex zeros of $F_\mu$ would require separate analysis.  No arithmetic
classification of $\mu$ follows here: the $A_j^\infty(r)$ are infinite
kernel values, not known rational numbers at a rational discount $r$.

\subsection{Removing the stub restriction from newborn comparison}

Let $W(s)$ again denote the supremum over all finite newborn environments,
including those with no inert stub.

\begin{theorem}[All newborn resolvents have one common critical scale]
\label{cr:all-newborns}
There is a finite constant $K_{\rm all}$ such that
\[
 W(s)\le K_{\rm all}V_s^\infty
 \qquad(\mu<s\le1).
\]
One may take $K_{\rm all}=1+20K/s_0^2$, where $K$ is from
Theorem~\ref{cr:uniform-inert}.  Consequently the bounded positive
critical eigenfunction construction extends to the entire newborn
class $\mathcal X$:
\[
 h(x_\infty)=1,\quad\gamma_1\le h\le K_{\rm all},
 \quad\mathcal K_\mu h=h.
\]
\end{theorem}
\begin{proof}
Cut the newborn tag's parent edge at time zero.  The retained state
starts at $x_\infty$.  In the actual process it has the inert generator
plus a predictable root-leaf input of intensity
\[
 q_t=1/D_p(t)\le1/2
\]
if the parent is real; the intensity is zero for an inert parent.
All outward selections from the retained component affect only the
exterior and are null in the retained state.  The exterior may respond
to the retained history; predictability is all that is needed below.

For the inert finite-horizon value $V_{s,r}(y)$, the root-leaf gradient
can be written with $W_T^\partial(s)$:
\[
 |\Delta_{\rm root}V_{s,r}(y)|
 \le\frac{40}{s}W_T^\partial(s)\qquad(0\le r\le T).
\]
Indeed, the two reference processes in this gradient retain their
original root stub.  Every charged newborn family is evaluated in
that full retained environment, so it belongs to
$\mathcal X^\partial$.  This is Lemma~\ref{cr:directional} with $r=0$.

Apply the finite-horizon discounted backward Dynkin identity to this
inert value along the actual controlled retained state.  It gives
\[
 V_{s,T}^{\rm tag}(x)=V_{s,T}^\infty+
 \mathbb E_x\int_0^T e^{-st}q_t
       \Delta_{\rm root}V_{s,T-t}(Y_t)\,dt.
\]
Thus, uniformly in the exterior,
\[
 V_{s,T}^{\rm tag}(x)
 \le V_{s,T}^\infty+\frac{20}{s^2}W_T^\partial(s).
\]
Take the supremum over $x$, let $T\to\infty$, and use
Theorem~\ref{cr:uniform-inert}.  This proves the displayed constant.
Finally normalize all newborn resolvents by $V_s^\infty$ and repeat
the diagonal and dominated-convergence proof of
Theorem~\ref{cr:eigenfunction}, now on $\mathcal X$.
\end{proof}

\subsection{A core-comparable critical generator eigenfunction}

Here the old state $y$ has exactly one inert stub, at its retained
root, and $V_s(y)$ counts all its old and future cores.  This restriction
is needed for the degree counting below.

\begin{lemma}[A positive density of old cores can found newborn families]
\label{cr:old-core-lower}
Put
\[
 p_* = \frac{e^{-2}(1-e^{-1})^2}{6},
 \qquad a_* = \frac{p_*e^{-2}}{192}>0.
\]
For every finite one-stub old state $y$ and $\mu<s\le1$,
\[
 V_s(y)\ge a_* C(y)V_s^\infty.
\]
\end{lemma}
\begin{proof}
Write $k=C(y)$ and let $c_v$ denote the number of core neighbors plus
the stub, where present.  Thus $c_v\ge1$ and
$\sum_vc_v=2k-1$.  Let $n_1$ be the number with $c_v=1$, let
$H=\{v:c_v\ge3\}$, and let
\[
 B=\{v:c_v=2,\ L_v=0\}.
\]
Every $c_v=1$ vertex has a leaf, since its full degree is at least two.
The degree sum implies $|H|\le n_1-1$.
Let $I\subset B$ consist of those vertices with no real neighbor in
$B$.  These form an independent set.  All its vertices have two real
core edges except possibly the root, which may use the unique stub.
Its incident core edges are distinct, so
$2|I|-1\le k-1$, or $|I|\le k/2$.

Call a core eligible if it has $c_v\le2$ and at least one leaf, or if
it belongs to $B\setminus I$.  For their number $e$ we have
\[
 e=k-|H|-|I|,\qquad e\ge n_1.
\]
Combining the inequalities gives $e\ge k/4$.

If an eligible core $v$ has a leaf, its first source-clock ring by
time one chooses a leaf with probability at least $1/3$: before that
ring, its core-neighbor count remains fixed at at most two and its
leaf count can only grow.  The event probability is at least
$(1-e^{-1})/3\ge p_*$.

Otherwise $v\in B$ has a neighbor $w\in B$.  During the first unit
interval keep $v$ and the other real neighbor of $w$, if any, silent.
This costs at most the factor $e^{-2}$.  Then $w$ keeps degree two and
no leaves during that interval.  Require its first ring to occur by
time one and select $v$, with probability $(1-e^{-1})/2$.  This gives
$v$ a leaf.  Further rings of $w$ are allowed.  During the second unit
interval the first ring of $v$ occurs with probability $1-e^{-1}$,
and selects a leaf with probability at least $1/3$.  Thus $v$ produces
a newborn core child by time two with probability at least $p_*$.
Other immigration and leaf births only improve the stated leaf-choice
lower bounds.  If $w$'s other neighbor is the stub, the unnecessary
silence condition can simply be omitted.

For each eligible old core retain at most the one child specified by
this event.  These newborns have mutually disjoint future genealogies,
even though their subsequent dynamics need not be independent.
Conditional on a child's birth at $\tau\le2$, its future discounted
mean is at least $\gamma_sV_s^\infty\ge V_s^\infty/48$ by the
existing newborn minorization.  Discounting the birth costs at most
$e^{-2s}\ge e^{-2}$.  Summing these nonnegative disjoint family rewards
over at least $k/4$ eligible old cores proves the assertion.
\end{proof}

\begin{theorem}[Critical eigenfunction in the core-comparable class]
\label{cr:generator-eigenfunction}
There is a function $H$ on the finite one-stub old states such that
\[
 H(x_\infty)=1,\qquad
 a_* C(y)\le H(y)\le A_*C(y),\qquad
 \mathcal L_\infty H=\mu H,
\]
for a finite constant $A_*$.  Moreover
\[
 P_t^\infty H(y)=e^{\mu t}H(y)
 \qquad(t\ge0).
\]
\end{theorem}
\begin{proof}
Normalize $H_s(y)=V_s(y)/V_s^\infty$ for $\mu<s\le1$.
Lemma~\ref{cr:old-core-lower} gives the uniform lower bound.
The old-core decomposition, using the closed stub subclass for each
future newborn family, gives
\[
 V_s(y)\le\frac{C(y)}s(1+W^\partial(s)).
\]
Since $V_s^\infty\ge1/s\ge1$, Theorem~\ref{cr:uniform-inert} yields
\[
 H_s(y)\le\frac{1+K}{s_0}C(y).
\]
Take a pointwise convergent diagonal subsequence on the countable old
state space.  At each fixed state the generator has only finitely many
possible distinct ordinary birth transitions.  Dividing the resolvent
equation by $V_s^\infty$ and passing to the limit therefore gives
\[
 (\mu-\mathcal L_\infty)H=0.
\]
Its normalization and both comparison bounds pass to the limit.

For the semigroup assertion, the finite-time resolvent identity is
\[
 V_s(y)=\int_0^t e^{-su}P_u^\infty C(y)\,du
       +e^{-st}P_t^\infty V_s(y).
\]
Divide by $V_s^\infty\to\infty$.  The first term tends to zero.
The uniform bound $H_s\le A_*C$ and finite-time Yule moments justify
dominated convergence in the second term, giving
$H(y)=e^{-\mu t}P_t^\infty H(y)$.
\end{proof}

In particular the elementary population-mode eigenfunctions that exist
at all positive parameters are no longer the only available harmonic
relations: the original critical parameter has an eigenfunction in the
specific core-comparable class.  This is an analytic selection theorem.
It still supplies no number-theoretic classification of that parameter.
Scaling $H$ by $a_*^{-1}$ gives an element of the class
$\mathcal A_C$ that attains the infimum in
Theorem~\ref{pm:variational}.

\begin{corollary}[The critical core-ring equation has no escape term]
\label{cr:core-ring-noescape}
For the one-stub eigenfunction $H$ of
Theorem~\ref{cr:generator-eigenfunction},
\[
 Q_\mu H=H.
\]
More precisely, if $T$ is the first core-clock ring from a state with
$k$ cores, then
\[
 \mathbb E_y[e^{-\mu t}H(Y_t);T>t]
 \le A_* k e^{-(\mu+k)t}\longrightarrow0.
\]
\end{corollary}
\begin{proof}
Before $T$, only leaf clocks ring, so the core count stays $k$, while
$\Pr(T>t)=e^{-kt}$.  Core comparison gives the displayed escape bound.
The semigroup identity makes $e^{-\mu t}H(Y_t)$ a true martingale.
Its stopped values at $T\wedge t$ are bounded by $A_*(k+1)$, since
the first core ring raises the core count by at most one.  Optional
stopping and dominated convergence therefore yield
$H(y)=\mathbb E_y[e^{-\mu T}H(Y_T)]=Q_\mu H(y)$.
\end{proof}

\begin{corollary}[Uniform exponential bounds and the free-tree eigenfunction]
\label{cr:free-eigenfunction}
The inert newborn mean satisfies
\[
 A_*^{-1}e^{\mu t}\le F_\infty(t)\le a_*^{-1}e^{\mu t}
 \qquad(t\ge0).
\]
Consequently, for every $s>\mu$,
\begin{equation}
 \frac1{A_*(s-\mu)}\le V_s^\infty
 \le\frac1{a_*(s-\mu)}.
 \label{cr:real-pole-order}
\end{equation}
This is a two-sided real-axis estimate; neither a meromorphic
continuation nor existence of a residue is asserted.
On the free old-state class with at least one core, there is likewise
a core-comparable function $H_{\rm free}$ with
\[
 \mathcal L_{\rm free}H_{\rm free}=\mu H_{\rm free},
 \qquad P_t^{\rm free}H_{\rm free}=e^{\mu t}H_{\rm free}.
\]
Consequently the original edge-start mean satisfies
$F_2(t)\asymp e^{\mu t}$ for $t\ge1$, with fixed positive comparison
constants.  Its embedded expected core count satisfies
$\mathbb E C_n\asymp n^\mu$ for $n\ge3$.
\end{corollary}
\begin{proof}
For the inert process, apply its positive semigroup to
$a_*C\le H\le A_*C$ and use $H(x_\infty)=1$.

For a free old state with $k\ge2$ cores, the same eligible-core proof
uses $\sum c_v=2k-2$, $|H|\le n_1-2$, and
$2|I|\le k-1$.  It therefore again gives at least $k/4$ eligible
cores and the same lower bound
$V_s^{\rm free}(y)\ge a_*kV_s^\infty$.
If $k=1$, the unique core has at least two leaves and no core neighbor;
its first source ring produces a child, so the same weaker constant
is valid directly.  The upper bound now uses
$W(s)\le K_{\rm all}V_s^\infty$ in the old-core decomposition.
Normalize by $V_s^\infty$ and repeat the generator and semigroup
limit proof.  This establishes $H_{\rm free}$ on the closed class
of states with at least one core.  No core-comparable function is
claimed at the initial edge itself, whose core count is zero.

The first free birth has independent exponential time of rate two,
and leaves the three-vertex star $S_3$.  Thus
\[
 F_2(t)=\int_0^t2e^{-2u}F_{S_3}(t-u)\,du.
\]
Core comparison for $H_{\rm free}$ gives
$F_{S_3}(t)\asymp e^{\mu t}$, and this convolution proves the assertion
for $t\ge1$.

If desired, the positive eigenfunction itself extends to the initial
edge $E_2$ by the uniquely required value
\[
 H_{\rm free}(E_2)=\frac{2H_{\rm free}(S_3)}{\mu+2}.
\]
The generator equation there is
$2(H_{\rm free}(S_3)-H_{\rm free}(E_2))
 =\mu H_{\rm free}(E_2)$, and the same first-birth decomposition
proves the semigroup equation at $E_2$.  Core comparison is still
asserted only when the state has a core.

Finally, on a free tree of population $n\ge2$ the generator equation
is exactly
\[
 \mathbb E[H_{\rm free}(T_{n+1})\mid T_n]
  =\left(1+\frac\mu n\right)H_{\rm free}(T_n).
\]
Starting from $E_2$, iteration gives the precise identity
\[
 \mathbb E H_{\rm free}(T_n)
 =H_{\rm free}(E_2)
   \prod_{j=2}^{n-1}\left(1+\frac\mu j\right)
 =H_{\rm free}(E_2)
   \frac{(\mu+2)_{n-2}}{(n-1)!}
 =H_{\rm free}(E_2)
   \frac{\Gamma(n+\mu)}{\Gamma(n)\Gamma(2+\mu)}.
\]
Its core-comparability and the elementary gamma-ratio asymptotic yield
$\mathbb E C_n\asymp n^\mu$.
\end{proof}

The last product also gives a nonnegative embedded-chain martingale
\[
 \frac{(n-1)!}{(\mu+2)_{n-2}}H_{\rm free}(T_n).
\]
The function $H_{\rm free}$ is selected by the limiting resolvents; an
explicit arithmetic formula for its values has not been proved.

% Consequence of the newly proved two-sided exponential mean bounds.
% This does not identify an arithmetic type or a limiting amplitude.

\subsection{A quantitative derivative-order bound at a fixed discount}

Let $H_{\rm free}$ be the critical core-comparable eigenfunction, and
write its comparison as
\[
 a_*C\le H_{\rm free}\le A_{\rm free}C
\]
on free states with a core.  One may use
$A_{\rm free}=(1+K_{\rm all})/s_0$ and the constant $a_*$ from
Lemma~\ref{cr:old-core-lower}.  The first-birth decomposition of the
edge-start process gives, for all $t\ge0$,
\begin{equation}
 a(e^{\mu t}-e^{-2t})\le F_2(t)
 \le b(e^{\mu t}-e^{-2t}),
 \qquad \frac ba=\frac{A_{\rm free}}{a_*}=:L,
 \label{cr:two-exponential-envelope}
\end{equation}
where
\[
 a=\frac{2H_{\rm free}(S_3)}{(\mu+2)A_{\rm free}},
 \qquad b=\frac{2H_{\rm free}(S_3)}{(\mu+2)a_*}.
\]

For a fixed $s>\mu$, retain
\[
 Z_m(s)=\int_0^\infty t^m e^{-st}F_2(t)\,dt,
 \qquad B_m(s)=Z_m(s)/m!.
\]
Put $\lambda=s-\mu$ and $\rho=\lambda/(s+2)\in(0,1)$.

\begin{proposition}[Derivative-ratio rate]
\label{cr:derivative-ratio-rate}
For every integer $m\ge0$,
\begin{equation}
 \left|\lambda\frac{B_{m+1}(s)}{B_m(s)}-1\right|
 \le\Delta_m(s):=
 \frac{L}{(1-\rho^{m+1})\sqrt{m+1}}.
 \label{cr:ratio-rate-forward}
\end{equation}
Whenever $\Delta_m(s)<1$,
\begin{equation}
 \left|\frac{B_m(s)}{B_{m+1}(s)}-(s-\mu)\right|
 \le (s-\mu)\frac{\Delta_m(s)}{1-\Delta_m(s)}.
 \label{cr:ratio-rate-inverse}
\end{equation}
In particular, at each fixed $s>\mu$,
\[
 \frac{B_m(s)}{B_{m+1}(s)}=(s-\mu)+O_s(m^{-1/2}).
\]
\end{proposition}
\begin{proof}
Set $A(t)=e^{-\mu t}F_2(t)$, and let $T$ be gamma distributed with
shape $r=m+1$ and rate $\lambda$.  By
\eqref{cr:two-exponential-envelope},
\[
 a(1-e^{-(\mu+2)t})\le A(t)\le b.
\]
Consequently
\[
 \mathbb EA(T)\ge a(1-\rho^r).
\]
Writing $U=\lambda T/r$, the gamma mean and variance give
$\mathbb EU=1$ and $\mathbb E(U-1)^2=1/r$.  Directly from the integral
definitions,
\[
 \lambda\frac{B_{m+1}(s)}{B_m(s)}
 =\frac{\mathbb E[A(T)U]}{\mathbb EA(T)}.
\]
Cauchy--Schwarz and $A\le b$ now show
\[
 \left|\frac{\mathbb E[A(T)(U-1)]}{\mathbb EA(T)}\right|
 \le\frac{b}{a(1-\rho^r)\sqrt r},
\]
which is \eqref{cr:ratio-rate-forward}.  The inverse-ratio bound follows
by taking the reciprocal of $1+\varepsilon_m$, with
$|\varepsilon_m|\le\Delta_m<1$.
\end{proof}

The constants in this bound are finite constants from the explicit
comparison constructions.  They have not been numerically evaluated
here, and the finite profile range used to prove their existence is
very large.  This establishes a derivative-order rate, not a practical
precision estimate for the existing calculator.  No monotonicity or
fixed one-sided sign of the ratio error is asserted.  Nor is a
power-law amplitude or a meromorphic singularity required.
For a certified lower bound $a_0\le\mu$ and $s>\mu$, replacing
$\rho$ by $(s-a_0)/(s+2)$ and $\lambda$ by $s-a_0$ gives an upper
error bound that does not require the unknown value of $\mu$.
A certified numerical bound on $L$ is still needed to use that
estimate with computed derivative intervals.
